\documentclass[10pt,a4paper]{amsart}
\usepackage[margin=19mm]{geometry}
\usepackage{amsmath,amssymb,mathtools}
\usepackage[scr=rsfs,cal=euler]{mathalfa}
\usepackage{xfrac}
\usepackage[unicode,hidelinks,bookmarks=false]{hyperref}
\newcommand{\ie}{i.e.\ }
\newcommand{\cf}{cf.\ }
\newcommand{\resp}{respectively\ }
\newcommand{\Subsection}{\subsection}
\providecommand{\nobreakdash}{\nobreak\hbox{-}}
\setlength{\emergencystretch}{2em}
\multlinegap0pt
\usepackage{bm}
\usepackage{bbm}
\usepackage{dsfont}
\usepackage{xspace}
\usepackage{cancel}
\usepackage{mathtools}
\usepackage{amsthm}
\makeatletter
\@ifundefined{claim*}{\newtheorem*{claim*}{Claim}}{}
\@ifundefined{theorem}{\newtheorem{theorem}{Theorem}}{}
\@ifundefined{corollary}{\newtheorem{corollary}[theorem]{Corollary}}{}
\@ifundefined{example}{\newtheorem{example}[theorem]{Example}}{}
\makeatother
\makeatletter
\newcommand{\Ad}{\operatorname{Ad}}
\newcommand{\Ann}{\operatorname{Ann}}
\newcommand{\Bl}{\operatorname{Bl}}
\newcommand{\Ima}{\operatorname{Im}}
\newcommand{\Z}{\mathbb{Z}}
\expandafter\ifx\csname customthm\endcsname\relax
\newenvironment{customthm}[1]
%  {\renewcommand\theinnercustomthm{#1}\innercustomthm}
\fi
\expandafter\ifx\csname proof\endcsname\relax
\renewenvironment{proof}[1][\relax]%s
%  {\paragraph{\textbf{Proof:}\ifx#1\relax\else~of #1\fi}}%
\fi
\expandafter\ifx\csname romanlist\endcsname\relax
\newenvironment{romanlist}
	{\begin{enumerate}
	\renewcommand{\theenumi}{\it\roman{enumi}}}
\fi
\newcommand{\that}{\hat{t}}
\makeatother
\title[Algebraic concordance of links]{Algebraic concordance of links}
\author{David Cimasoni, Anthony Conway, Gaetan Simian}
\date{Source: arXiv:2607.25972v1 (CC BY 4.0)}
\begin{document}
\maketitle

\noindent\emph{Source and licence.} Algebraic concordance of links. Version arXiv:2607.25972v1 (CC BY 4.0), distributed under the Creative Commons Attribution 4.0 International licence, as recorded in the arXiv licence metadata for this exact version and shown on the abstract page https://arxiv.org/abs/2607.25972v1. Full rights notice: licenses/arxiv-2607.25972-CC-BY-4.0.txt.\par\smallskip

\noindent\textit{Editorial scope.} Counted unit: the Theorem whose source label is \texttt{theorem:Invariance3ManifoldVersion} in the article (source0.tex lines 981--987, source label \texttt{theorem:Invariance3ManifoldVersion}), delivered together with the article's own complete original proof of it (source0.tex lines 988--1060, one \texttt{proof} environment of 73 lines). No mathematics is written, repaired or re-derived: statement and proof are the article's own text line for line. Required context retained uncounted: cor:max-pn (lines 4086--4090); ex:NoPN (lines 4095--4102). Every definition, display and label of those lines is retained unchanged. Omitted from this package and not redistributed: 1--980, 1061--4085, 4091--4094, 4103--4176. Nothing inside a retained excerpt is omitted or altered: every retained body is delivered exactly as the article's lines are written, so no omission receipt of any kind accompanies this package. Citations: the retained excerpts carry no citation command at all, so no bibliography entry of the article is transcribed and no reference is redistributed. Reference adaptation: none was needed and none was made; every reference inside the delivered unit resolves inside it, so no unresolved reference or citation is produced. Printed slips kept literal: no printed word, symbol, hypothesis or number of the counted statement or of its proof was altered. Layout and class: the article's own class is \texttt{amsart} with packages including a4wide, babel, ulem, natbib, xy, pstricks, enumerate, amsfonts, amssymb, amsmath, pinlabel, array, hhline, slashed; the wrapper is the lane's pinned \texttt{amsart} editorial wrapper at 10pt on A4, which restates the theorem-like environments the retained text needs and adds the article's own macro definitions that the retained text needs (\texttt{\textbackslash Ad}, \texttt{\textbackslash Ann}, \texttt{\textbackslash Bl}, \texttt{\textbackslash Ima}, \texttt{\textbackslash Z}, \texttt{\textbackslash that}), with extra wrapper packages (bm, bbm, dsfont, xspace, cancel, mathtools). The delivered numbering is the unit's own. Assets: no retained span contains a figure environment, a graphics-inclusion command, a PDF-page inclusion, a table or a TikZ picture; no image file is copied into this package, the package declares no asset dependency and no image file is redistributed with it. Licence: version arXiv:2607.25972v1 of this article is distributed under the Creative Commons Attribution 4.0 International licence, as recorded in the arXiv licence metadata for this exact version and shown on the abstract page https://arxiv.org/abs/2607.25972v1. A full rights notice is delivered at licenses/arxiv-2607.25972-CC-BY-4.0.txt. Characters: every retained excerpt is pure ASCII, so the delivered engine, XeLaTeX with its default Latin Modern text font, reports no missing character.
\begin{corollary}
\label{cor:max-pn}
If~$R$ is a UFD, then the maximal pseudonull submodule~$zM$ of an~$R$-module~$M$ is generated by the
elements of~$TM$ whose annihilator ideal contains two coprime elements.\qed
\end{corollary}

\begin{example}
\label{ex:NoPN}
If~$R$ is a UFD, then the~$R$-module~$R_0/R$ has trivial maximal pseudonull submodule.
To show this, consider~$m\in R_0/R$ admitting two coprime elements~$\alpha,\beta\in\Ann(m)$. We write~$m\in R_0/R$ as the class of a fraction~$\frac{r}{s}\in R_0$ with~$r,s \in R$ coprime and~$s\neq 0$. By assumption, both~$\frac{\alpha r}{s}$ and~$\frac{\beta r}{s}$ are elements of~$R$, so~$s$ divides~$\alpha r$ and~$\beta r$ in~$R$.
Hence, any prime factor~$p$ of~$s$ divides both~$\alpha r$ and~$\beta r$ in~$R$; since~$r$ and~$s$ are coprime,~$p$ divides the coprime elements~$\alpha$ and~$\beta$, a contradiction.
Therefore, the element~$s\in R$ is a unit, showing that~$m=0\in R_0/R$.
By Corollary~\ref{cor:max-pn}, this implies that the maximal pseudonull submodule of the~$R$-module~$R/R_0$ is trivial.
\end{example}

\begin{theorem}
\label{theorem:Invariance3ManifoldVersion}
Let~$Y$ be a closed~$3$-manifold and let~$\varphi \colon H_1(Y) \to \Z^\mu$ be an epimorphism.
If~$(Y,\varphi)$ bounds a~$4$-manifold~$(Z,\varphi)$ such that the sequence~$
TH_2(Z, \partial Z; \Lambda) \stackrel{\partial}{\longrightarrow} TH_1(\partial Z; \Lambda) \stackrel{\iota}{\longrightarrow} TH_1(Z; \Lambda)$ is exact,
then~$\hat \Bl_Y^\varphi$ is metabolic.
\end{theorem}

\begin{proof}
During this proof,  we will omit the~$\varphi$ superscripts for brevity.
The strategy of the proof is to show that~$\Bl_{\partial Z}=\Bl_Y$ is metabolic and to then deduce,  by a purely algebraic argument,  that~$\hat \Bl_{\partial Z}$ is metabolic.
During the proof of this first step,  we will use the hypothesis that the exact sequence of the pair gives rise to the exact sequence~$TH_2(Z, \partial Z; \Lambda) \stackrel{\partial}{\longrightarrow} TH_1(\partial Z; \Lambda) \stackrel{\iota}{\longrightarrow} TH_1(Z; \Lambda)$.
The following claim carries out the first step of our plan.
\begin{claim*}
For~$K:=\ker(\iota)=\Ima(\partial)$, the following is a metabolizer for~$\Bl_{\partial Z}$:
\[
P:=K^\perp=\{ x\in TH_1(\partial Z;\Lambda)\mid \Bl_{\partial Z}(x,y)=0 \text{ for all } y\in K\}. 
\]
\end{claim*}
\noindent To prove this claim, consider the Blanchfield pairing
$
\Bl_Z \colon TH_2(Z,\partial Z; \Lambda) \times TH_1(Z;\Lambda) \to Q/\Lambda.
$
A now standard diagram chase yields the inclusion $K \subset K^\perp$, which implies~$P^\perp=K^{\perp\perp}\subset K^\perp=P$.
%%Don't delete
%The adjoints for the pairings~$\Bl_{\partial Z}$ and~$\Bl_Z$ fit in the following diagram which commutes by naturality of Poincar\'e duality, of the Bockstein homomorphism, and of the evaluation map:
%\begin{center}
%\begin{tikzcd}
%TH_1(\partial Z;\Lambda) \arrow[bend left = 15]{rrr}{\Ad_{\Bl_{\partial Z}}} \arrow{d}{\iota} \arrow{r}{\PD^{-1}} & TH^2(\partial Z;\Lambda) \arrow{r}{\BS^{-1}} \arrow{d}{\delta} & \frac{H^1(\partial Z;Q/\Lambda)}{\ker(\BS)} \arrow{r}{\ev} \arrow{d}{\delta} & \overline{\Hom(TH_1(\partial Z;\Lambda),Q/\Lambda)}\arrow{d}{\partial^*} \\
%TH_1(Z;\Lambda) \arrow[bend right =15, swap]{rrr}{\Ad_{\Bl_Z}} \arrow{r}{\PD^{-1}} & TH^3(Z,\partial Z;\Lambda)  \arrow{r}{\BS^{-1}} & \frac{H^2(Z,\partial Z;Q/\Lambda)}{\ker(\BS)} \arrow{r}{\ev} & \overline{\Hom(TH_2(Z,\partial Z;\Lambda),Q/\Lambda)}\,.\\
%\end{tikzcd}
%\end{center}
%%which commutes by naturality of Poincar\'e duality, of the Bockstein homomorphism, and of the evaluation map.
%Thus, for any~$X\in TH_2(Z,\partial Z; \Lambda)$ and any~$y \in TH_1(\partial Z; \Lambda)$, we have
%\begin{equation}
%\label{eq:BlanchfieldCommute}
%\begin{split}
%\Bl_Z(X, \iota(y)) =& \Ad_{\Bl_Z}(\iota(y))(X) =\left((\Ad_{\Bl_Z}\circ \iota)(y)\right)(X)=\left((\partial^\ast\circ\Ad_{\Bl_{\partial Z}})(y)\right)(X)\\=&\left(\Ad_{\Bl_{\partial Z}}(y)\right)(\partial X)=\Bl_{\partial Z}(\partial X, y)\,.
%\end{split}
%\end{equation}
%For all~$x\in K=\Ima(\partial)$ and~$y\in K=\ker(\iota)$, there exists~$X\in TH_2(Z,\partial Z;\Lambda)$ such that~$\partial X=x$ and we have
%\[
%\Bl_{\partial Z}(x,y)=\Bl_{\partial Z}(\partial X,y)\stackrel{\eqref{eq:BlanchfieldCommute}}{=}\Bl_Z(X, \iota(y))=\Bl_Z(X,0)=0\,.
%\]
%%%

We now check the opposite inclusion~$K^{\perp}\subset K^{\perp\perp}$,
including the details to verify that the argument carries over from~$\Lambda_S$ to~$\Lambda$.
To do so, first observe that for any~$y\in K^\perp$, its image~$\iota(y)\in\iota(K^\perp)$ satisfies~$\Bl_Z(X,\iota(y))=\Bl_{\partial Z}(\partial X,y)=0$
for all~$X\in TH_2(Z,\partial Z;\Lambda)$, since~$\partial X$ belongs to~$K=\Ima(\partial)$ and~$y$ to~$K^\perp$.
Hence,~$\iota(y)$ lies in the kernel of~$\Ad_{\Bl_{Z}}$, which by Blanchfield's theorem is the maximal pseudonull submodule of~$TH_1(Z;\Lambda)$. Since we know that the inclusion~$K\subset K^\perp$ holds, this implies
\[
K^\perp/K=K^\perp/\ker(\iota)\cong\iota(K^\perp)\subset zH_1(Z;\Lambda)\,.
\]
Assume that~$x$ belongs to~$K^\perp$. By the equation displayed above together with Corollary~\ref{cor:max-pn},
the class~$[x]\in K^\perp/K$ can be written as~$[x]=\sum_i[x_i]$ where for each~$i$, there exists coprime
elements~$\alpha_i,\beta_i\in\Lambda$ such that~$\alpha_i[x_i]=\beta_i[x_i]=0\in  K^\perp/K$; in other words,
the elements~$\alpha_i x_i$ and~$\beta_i x_i$ belong to~$K$. Hence, for all~$i$ and all~$y\in K^\perp$, we have~$\alpha_i\Bl_{\partial Z}(x_i,y)=\Bl_{\partial Z}(\alpha_ix_i,y)=0$
and similarly~$\beta_i\Bl_{\partial Z}(x_i,y)=0$. Since~$\alpha_i$ and~$\beta_i$ are coprime, Corollary~\ref{cor:max-pn}
and Example~\ref{ex:NoPN} imply~$\Bl_{\partial Z}(x_i,y)=0$ for all~$i$ and all~$y\in K^\perp$.
This yields~$\Bl_{\partial Z}(x,y)=0$ for all~$y\in K^\perp$, meaning that~$x$ belongs to~$K^{\perp\perp}$.
This shows that~$P=K^\perp=K^{\perp\perp}=P^\perp$, thus proving the claim.


Let~$\that\colon TH_1(\partial Z;\Lambda)\to \that H_1(\partial Z;\Lambda)$ denote the canonical projection.
We now prove that~$\that(P)$ is a metabolizer for~$\hat \Bl_{\partial Z}$.
By definition, note that~$\hat \Bl_{\partial Z}(\that(x),\that(y)) = \Bl_{\partial Z}(x,y)$
for all~$x,y\in TH_1(\partial Z;\Lambda)$.
For any submodule~$N\subset TH_1(\partial Z;\Lambda)$, this equation readily implies 
\begin{equation}
\label{eq:hat-perp}
\that(N)^\perp=\that(N^\perp).
\end{equation}
The claim establishes that the submodule~$P=K^\perp\subset TH_1(\partial Z;\Lambda)$ satisfies~$P^\perp=P$. By~\eqref{eq:hat-perp} applied to~$N=P$, its image~$\that(P)\subset \that H_1(\partial Z;\Lambda)$
satisfies
\[
\that(P)^\perp=\that(P^\perp)=\that(P)\,.
\]
In other words, the submodule~$\that(P)\subset \that H_1(\partial Z;\Lambda)$ is a metabolizer for the nondegenerate pairing~$\hat \Bl_{\partial Z}$. 
This concludes the proof.
\end{proof}

\end{document}
